\reportfigure{\begin{tikzpicture}
\begin{axis}[width=.94\linewidth,height=5.7cm,xmin=-340,xmax=340,ymin=0,ymax=1350,ytick={0,300,600,900,1200},xtick={-300,-150,0,150,300},
xlabel={Distance from the fire axis (cm)},ylabel={Temperature rise (K)},axis lines=left,ymajorgrids,grid style={gray!15},
legend columns=3,legend style={at={(.5,1.03)},anchor=south,draw=none,/tikz/every even column/.append style={column sep=6pt}}]
\addplot[gray,thick,dashed] table[col sep=comma,x expr=\thisrow{model_x_cm}-330,y=fds_rise_K]{../../05_Calculations/op6_106_doorflow_90_120.csv};
\addlegendentry{Adiabatic}
\addplot[teal,thick] table[col sep=comma,x expr=\thisrow{model_x_cm}-330,y=fds_rise_K]{../../05_Calculations/op6_106_board_doorflow_90_120.csv};
\addlegendentry{Conducting board}
\addplot[only marks,rust,mark=o,mark size=1.6pt] table[col sep=comma,x=figure_r_cm,y=exp_rise_K]{../../05_Calculations/op6_106_board_doorflow_90_120.csv};
\addlegendentry{Peng et al.}
\end{axis}\end{tikzpicture}}{Figure 2. OP-6 at 106.2 kW. The adiabatic and conducting-board predictions share the doorway diagnostic grid. The measured and simulated profiles use the same coordinate origin as Figure 1.}
